\section{Sixfold cycles after Schoen and Koike, and bounds for every fourfold determinant}
\label{sec:sixfolds}
This section keeps track of the degrees in Schoen's cyclic-cover construction and improves the passage between integral polarization lattices.

Write $U=\left(\begin{smallmatrix}0&1\\1&0\end{smallmatrix}\right)$ for a Hermitian hyperbolic plane. Degrees are unnormalized intersection numbers. All polarizations below are compatible with the indicated quadratic action.

\begin{theorem}\label{thm:sixfold-suppliers}
For full maximal-order action, the following detector bounds hold on the indicated principal split sixfold families:
\[
\begin{array}{lll}
K&\text{Hermitian lattice}&\delta_W\text{ at most}\\
\hline
\Q(i)&U^3&746496\\
\Q(i)&I_{3,3}&5971968\\
\Q(\omega)&U^3\cong I_{3,3}&30375.
\end{array}
\]
For an arbitrary full maximal-order split sixfold of polarization product $\Delta$, one has
\[
 \delta_W\le5971968\Delta\quad(K=\Q(i)),
 \qquad \delta_W\le30375\Delta\quad(K=\Q(\omega)).
\]
For an even Gaussian Hermitian lattice the first bound improves to $746496\Delta$.
\end{theorem}

\begin{theorem}\label{thm:all-determinant-fourfolds}
Let $(A,L)$ be a Weil fourfold over $K$, of any determinant, and put $\Delta=\chi(L)$. Suppose the order $\mathcal O_f=\Z+f\mathcal O_K$ acts integrally. Then
\[
\begin{aligned}
 \delta_W(A,L)&\le1990656\,f^6\Delta^2
 &&\text{if }K=\Q(i),\\
 \delta_W(A,L)&\le10125\,f^6\Delta^2
 &&\text{if }K=\Q(\omega).
\end{aligned}
\]
Each right-hand side may be divided by $d$, where
$d=[\mathcal O_KH_1(A,\Z):H_1(A,\Z)]$.
\end{theorem}

The geometric inputs are published: Schoen constructs the cycles \cite[Theorem~2.0 and \S3]{Schoen88}, and Schoen and Koike prove the required dominance \cite[Lemma~13, pp.~334--335]{Schoen98}, \cite[Theorem~2.1, p.~568]{Koike}. The estimates below follow from their constructions, integral polarization calculations and closedness of bounded detector loci. They do not use the recent CM theorems or equivariant semiregularity.

\subsection{The cyclic sources and their principal polarizations}

Let $C\to X$ be a connected cyclic \'etale cover of degree $m=3$ or $4$, with $g(X)=4$ and generator $\sigma$. Then $g(C)=3m+1$. Choose a topological symplectic basis on $X$ for which all monodromy lies on the first handle. The other three handles lift in $m$ disjoint copies. Their cycles satisfy
\[
 E(a_{j,r},b_{k,s})=\delta_{jk}\delta_{rs},
 \qquad E(a_{j,r},a_{k,s})=E(b_{j,r},b_{k,s})=0,
\]
and $\sigma$ shifts the sheet index. The cover of the first punctured handle is a genus-one surface with $m$ boundary circles. After the other handles are glued, its boundary classes vanish and it contributes an invariant rank-two symplectic summand. This gives an integral orthogonal decomposition into that summand and three pairs of regular representations.

For $m=4$, the saturated primitive lattice on each handle is generated by
\[
 a_0-a_2,\quad a_1-a_3,\quad b_0-b_2,\quad b_1-b_3.
\]
The two nonzero alternating pairings are both two. Thus on the primitive sixfold $B=\ker(1+\sigma^2)^0$ the restriction of the Jacobian polarization is $2\Xi$, where $\Xi$ is principal. The Gaussian Hermitian matrix for $\Xi$ on a handle has off-diagonal entries $i,-i$; changing a basis vector by a unit gives $U$. The principal lattice is $U^3$.

For $m=3$, the norm Prym $P$ has, on each handle, augmentation generators $A=a_0-a_1$ and $B=b_0-b_1$ over $\Z[\omega]$. The form
$h(x,y)=E(x,\omega y)-\bar\omega E(x,y)$ has
\[
 E(A,B)=2,\quad E(A,\omega B)=-1,
 \qquad h(A,B)=1+2\omega=\sqrt{-3}.
\]
Its matrix is $\left(\begin{smallmatrix}0&\sqrt{-3}\\-\sqrt{-3}&0\end{smallmatrix}\right)$ and its alternating type is $(1,3)$. Enlarge the lattice on each handle from $\langle A,B\rangle$ to $\langle A,B/\sqrt{-3}\rangle$. This has index three and gives matrix $U$. On three handles it defines a degree-$27$ quotient $P\to P'$ with a principal polarization $\Xi'$ satisfying
\[
 \varphi^*\Xi'=\Theta_C|_P.
\]
The positive Riemann form on the real vector space has not changed. A finite level cover chooses the kernels in families and identifies the period domains. The Eisenstein basis with columns $(1,-\omega)$ and $(1,\bar\omega)$ takes $U$ to $\operatorname{diag}(1,-1)$ and has unit determinant. Thus the principal lattice is also $I_{3,3}$.

Koike's Theorem~2.1 states that the quartic Prym map is dominant onto the nine-dimensional unitary domain; its codifferential is multiplication of inverse primitive eigendifferentials. Schoen's Lemma~13 proves the cubic dominance by exhibiting nine independent products. These statements apply to the principalized sources just described, since a fixed isogeny change of lattice does not change the period domain.

\subsection{An effective detecting component and its degree}

Put
\[
 N=\ker\bigl((\Z/m)^6\xrightarrow{\sum}\Z/m\bigr),
 \qquad G_0=N\rtimes S_6,
 \qquad W_0=C^6/G_0.
\]
The finite maps
\[
 C^6\xrightarrow{r}C^{(6)}\xrightarrow{p}W_0
 \xrightarrow{\beta}X^{(6)}
\]
have generic degrees $6!,m^5,m$. Schoen's local equation \cite[Lemma~2.2, p.~13]{Schoen88} makes $\beta$ \'etale in the unramified case. The inverse image of $|K_X|\cong\mathbf P^3$ is a disjoint union of $m$ copies of $\mathbf P^3$. Choose one, $Q$.

The nonvanishing is especially simple here. Embed $X^{(6)}$ in $X^{(7)}$ by adding a point. Since $X^{(7)}\to\operatorname{Pic}^7(X)$ is a projective bundle, the canonical $\mathbf P^3$, after this embedding, is a fibre. Its normal bundle there is trivial of rank four; the normal line of $X^{(6)}$ restricts to $\mathcal O_{\mathbf P^3}(1)$. The normal sequence gives
\[
 c(N_{\mathbf P^3/X^{(6)}})=(1+H)^{-1},
 \qquad c_3(N_{\mathbf P^3/X^{(6)}})=-H^3\ne0.
\]
The same holds for $Q\subset W_0$. For a primitive character $\chi$ of the residual deck group, the cycle
$\sum_t\chi(-t)\,\sigma^tQ$ restricts on $Q$ to this nonzero top Chern class, since the other sheets are disjoint. Thus the primitive-character projection is nonzero. This is the unramified calculation of \cite[Lemma~2.6, p.~16]{Schoen88}.

Let $Z$ be an irreducible component of the inverse image of $Q$ in $C^6$ and let $W=r(Z)$ with reduced structure. A general canonical divisor of $X$ has six distinct zeros, so the relevant action is free at the generic point above it. If $H_Z\subset G_0$ stabilizes $Z$, then
\[
 \deg(Z/Q)=|H_Z|,
 \quad \deg(Z/W)=|H_Z\cap S_6|,
 \quad e:=\deg(W/Q)\le m^5.
\]
The primitive determinant plane $U_C\subset H^1(C)^{\otimes6}$ is $G_0$-invariant: the sum-zero sheet translations act trivially, and the alternating determinant and the cohomological permutation signs cancel. Its primitive-character components are those just detected on $W_0$. For $q:C^6\to W_0$, the generic cycle identity is
\[
 \sum_{g\in G_0}g_*[Z]=\deg(Z/Q)q^*[Q].
\]
Pair with an element of $U_C$ for which the right side is nonzero. All summands on the left have the same pairing, so $Z$ itself detects. This produces an effective integral component, without interpreting cyclotomic coefficients as effective cycles.

Let $\xi$ be the incidence class on $C^{(6)}$, and let $\alpha:C^{(6)}\to J(C)$ be an Abel map. Pullback of an incidence hyperplane of $|K_X|$ has class $m\xi$. Consequently
\begin{equation}\label{eq:cyclic-incidence}
 m^3\int_W\xi^3=e,\qquad \int_W\xi^3\le m^2.
\end{equation}
Write $\vartheta=\alpha^*\Theta_C$ and let $\mathcal D$ be the reduced big diagonal of $C^{(6)}$. Then
\begin{equation}\label{eq:symmetric-theta}
 \vartheta=(g(C)+5)\xi-\tfrac12[\mathcal D].
\end{equation}
To verify this, pull back to $C^6$. The incidence class becomes the sum of the point classes. The pair diagonal is the sum of two point classes minus the degree-one cross pairing; the Abel sum pulls theta back to $g(C)$ times the sum of point classes plus these cross pairings. Sum over the fifteen pairs and use $r^*\mathcal D=2\sum_{i<j}\Delta_{ij}$. Rational pullback is injective for the finite quotient.

The class $\xi$ is ample because its finite pullback is ample, and $\vartheta$ is nef. Since $W$ is not contained in $\mathcal D$, the restricted diagonal is effective. With $b=g(C)+5$,
\[
 (b\xi)^3-\vartheta^3
 =\tfrac12[\mathcal D](b^2\xi^2+b\xi\vartheta+\vartheta^2)
\]
has nonnegative intersection with $W$, computed on a resolution if necessary. This uses effectivity of the restricted diagonal, not its nefness. Hence
\begin{equation}\label{eq:cyclic-theta-degree}
 \int_W\vartheta^3\le b^3m^2.
\end{equation}

For $m=4$ use $\phi=(1-\sigma^2)\alpha$ with image in $B$. Its pullback of the Weil plane to $C^6$ is $U_C$, up to nonzero scalars. Indeed expand the exterior product of six primitive one-forms under the Abel sum. Two forms of the same primitive character on one factor multiply to zero, since their product has nontrivial character whereas $H^2(C)$ has trivial character. The surviving alternating tensor is nonzero; $1-\sigma^2$ is invertible on the primitive eigenspaces. Projection formula shows that $\phi(W)$ has dimension three and detects.

The theorem of the square and $\Theta_C|_B=2\Xi$ give
\[
 (1-\sigma^2)^*\Xi
 =2\Theta_C-\tfrac12(1+\sigma^2)^*\Theta_C
 \le2\Theta_C.
\]
Thus its effective pushforward has degree at most
\[
 2^3\cdot18^3\cdot4^2=746496.
\]
For $m=3$, $\alpha(W)$ lies in a translate of the norm Prym. The same argument using the invertible endomorphism $1-\sigma$ proves detection, and its restricted-theta degree is at most
\[
 15^3\cdot3^2=30375.
\]
The principal quotient has pullback polarization equal to that restricted theta, so adds no degree factor. Selecting a reduced image can only decrease either bound.

The bounds are uniform on the dense cyclic-source loci. Bounded-detector closedness extends them to the whole principal components. For the odd Gaussian bound, the three-block matrix with blocks $M=\left(\begin{smallmatrix}1&1\\1&-1\end{smallmatrix}\right)$ satisfies $M^*I_{1,1}M=2U$. It gives an isogeny from the even principal family to the odd one with polarization multiplier two. Threefold degrees increase by at most eight, giving $5971968$.

These are all Gaussian principal parity types. Split off a primitive isotropic vector and a vector pairing to one in a unimodular lattice. The partner norm reduces modulo two. An even norm gives $U$ after subtracting an integral multiple of the isotropic vector; an odd norm gives $I_{1,1}$. Repeat. A mixture containing an odd block becomes $I_{3,3}$: in $\langle1\rangle\oplus U$, with basis $e,a,b$, the vectors $e+a,e-b,e+a-b$ are a unimodular orthogonal basis of norms $1,1,-1$. This converts each remaining hyperbolic block. The rational isotropic spaces exist because the Hermitian space is split.

Theorem~\ref{thm:principal-quotient}, with $n=3$, transfers the principal bounds to every split polarization: $t_K^3\Delta$ is $8\Delta$ for the Gaussian ring and $\Delta$ for the Eisenstein ring. In the even Gaussian case the final paragraph of Section~\ref{sec:principal-quotients} gives the factor $\Delta$. This finishes Theorem~\ref{thm:sixfold-suppliers}.

\subsection{All determinants and nonmaximal orders}

\begin{proof}[Proof of Theorem~\ref{thm:all-determinant-fourfolds}]
First suppose the full maximal order acts. The balanced surface of Theorem~\ref{thm:explicit-stabilization} makes $A\times E_K^2$ a split sixfold of polarization product $\Delta^2$. In both fields $\kappa_K(\Delta)=1$, so the sixfold bound and the factor $1/3$ give
\[
 \delta_W(A,L)\le\frac{5971968\Delta^2}{3}
 =1990656\Delta^2
\]
in the Gaussian case, and $30375\Delta^2/3=10125\Delta^2$ in the Eisenstein case. The Gaussian product contains a norm-one vector from the surface, so no even-lattice improvement is used.

For conductor $f$, put $\Lambda'=\mathcal O_K\Lambda=\Lambda+\tau\Lambda$ and $d=[\Lambda':\Lambda]$. The form $fE$ is integral on $\Lambda'$: cross terms use the integral endomorphism $f\tau$, while
$E(\tau x,\tau y)=\Nm(\tau)E(x,y)$ follows from Rosati compatibility. The quotient isogeny $q:A\to A'$ has a compatible polarization with
\[
 q^*\theta'=f\theta,
 \qquad \Delta'=f^4\Delta/d.
\]
Apply the maximal-order bound $C(\Delta')^2$, where $C$ is the appropriate constant. Finite flat pullback preserves integral effectivity and a nonzero Weil projection, and multiplies the surface degree by $d/f^2$. Selecting a detecting irreducible component gives
\[
 \delta_W(A,L)\le\frac d{f^2}C\left(\frac{f^4\Delta}{d}\right)^2
 =C\frac{f^6}{d}\Delta^2.
\]
\end{proof}

The principal quotient improves transfer once a principal detector is known; it supplies no principal detector in a new dimension. Cyclic \'etale covers of genus-five curves have twelve moduli, whereas the signature-$(4,4)$ unitary domain has dimension sixteen. The genus-four dominance used here therefore does not extend to eightfolds by the same unramified cyclic construction.
